\scanpage{435}
\mspage{2}
% archive sheet number: 434

Si

\[
f\colon S\unsure{\mathrel{-}|\longrightarrow}T
\]

est un foncteur, \deleted{\unsure{tout}} on va définir un foncteur

\[
f_!\colon\uline{\mathrm{Grp}}(S)\longrightarrow\uline{\mathrm{Grp}}(T)
\]

(image directe de 1-types d'homotopie relatifs) Mais notons d'abord que si

\[
p\colon X\longrightarrow S
\]

est une $S$-catégorie \uline{quelconque} (pas néc. Cat-fibrée), on \unsure{peut associer|lui associe} néanmoins un type d'hom relatif, via la catégorie fibrée \inserted{$\widetilde X^{\unsure{(S)|S}}$ ou $\widetilde X$} ayant comme fibres les ${}_{s}\backslash X$ (en notant que $s'\longrightarrow s$ dans $S$ définit bien ${}_{s}\backslash X\longrightarrow {}_{s'}\backslash X$). \deleted{\unsure{Si $X$ est déjà}} On a un $S$-foncteur canonique

\margin{Ayant $\widetilde X$, on \unsure{peut définir|définit} $\Pi(X/S)\overset{\text{déf}}{=}\Pi(\widetilde X/S)$.}
% slanted

\[
X\xrightarrow{i}\widetilde X\qquad
x\longmapsto\bigl(x,p(x),\underset{\text{``}s\text{''}}{\deleted{\unsure{p'}}p(x)}\xrightarrow{\mathrm{id}}p(x)\bigr).
\]

\deleted{qui} C'est une $S$-\inserted{Hot-}équivalence si $X$ \unsure{était} déjà fibrée. -- En fait, en chaque fibre c'est un \deleted{\unsure{équivalé}} \unsure{homotopisme}, car on construit alors $\widetilde X_{s_0}\xrightarrow{j}X_{s_0}$ de façon évidente, de façon que

\diagram{\begin{tikzcd}[ampersand replacement=\&,column sep=12mm]
X_{s_0}\arrow[r,shift left,dash,"i"] \& \widetilde X_{s_0}\arrow[l,shift left,"j"]
\end{tikzcd}}

$ji\simeq\mathrm{id}$, $ij\xrightarrow[\text{\unsure{simple}}]{\text{homotopie}}\mathrm{id}$ (NB cette flèche n'est pas iso en général.)

Revenons alors à $f\colon S\longrightarrow T$, et considérons un groupoïde fibré sur $S$

\[
G\longrightarrow S,
\]

ce n'est pas une catégorie fibrée sur $\deleted{\unsure{S}}T$, \deleted{\unsure{p}} mais on peut lui associer

\[
t\longmapsto {}_{t}\backslash G\qquad T^{\unsure{\circ|o}}\unsure{\mathrel{-}|\longrightarrow}(\mathrm{Cat}),
\]

d'où une catégorie fibrée sur $T$, dont on prend le $\Pi$ relatif. Mais il est plus simple de prendre \deleted{les} \inserted{\unsure{directmt|directement}}

\[
t\longmapsto\Pi({}_{t}\backslash G)
\]

qui \deleted{\unsure{est}} \inserted{définit} une \deleted{\unsure{champ}} catégorie fibrée en groupoïdes sur $S$.
